\begin{unverified}
Status: agent-written proof note for the coordinate canonization search. It is
not reviewed by Jörn. The theorem is an exact-arithmetic generic statement.
The Rust implementation is an f64 prototype and only approximates this
construction.
\end{unverified}

\subsection{Generic coordinate canonization by a symplectic frame}
\label{sec:generic-coordinate-canonization}

This note isolates the part of the canonization problem that remains after the
polytope has already been scale-fixed and centered. The row covectors are
represented as column vectors in the fixed coordinate convention
\((q_1,q_2,p_1,p_2)\). Let \(V=\mathbb R^4\) with standard symplectic matrix
\(J\), and write
\[
  \omega(u,v)=u^T Jv
\]
for this coordinate-level pairing of represented covectors. A row list is
\[
  A=(a_1,\ldots,a_F)\in (V^*)^F.
\]

\begin{definition}[Omega row signatures]
\label{def:omega-row-signatures}
For a row list \(A\), define the skew matrix
\[
  \Omega(A)_{ij}=\omega(a_i,a_j).
\]
The row signature of index \(i\) is the sorted tuple
\[
  \sigma_i(A)=\operatorname{sort}\{\Omega(A)_{ij}:1\le j\le F\}.
\]
The sort order is fixed once and for all.
\end{definition}

\begin{definition}[Polynomial omega row signatures]
\label{def:polynomial-omega-row-signatures}
For \(1\le k\le F\), define
\[
  p_{i,k}(A)=\sum_{j=1}^F \Omega(A)_{ij}^k.
\]
Let
\[
  p_i(A)=(p_{i,1}(A),\ldots,p_{i,F}(A)).
\]
The vector \(p_i(A)\) is symmetric in the entries of row \(i\) of
\(\Omega(A)\). Therefore it depends only on the multiset underlying
\(\sigma_i(A)\).
\end{definition}

\begin{definition}[Generic row-labeling domain]
\label{def:generic-row-labeling-domain}
Let \(U_{\mathrm{label}}\subset (V^*)^F\) be the set of row lists such that the
tuples \(\sigma_1(A),\ldots,\sigma_F(A)\) are pairwise distinct. For
\(A\in U_{\mathrm{label}}\), let
\[
  \pi_A=(i_1,\ldots,i_F)
\]
be the unique ordering satisfying
\[
  \sigma_{i_1}(A)<\sigma_{i_2}(A)<\cdots<\sigma_{i_F}(A)
\]
in lexicographic tuple order. Define the ordered list
\[
  B_A=(b_1,\ldots,b_F)=(a_{i_1},\ldots,a_{i_F}).
\]
\end{definition}

\begin{lemma}[The row-labeling condition is open dense]
\label{lem:generic-coordinate-labeling-open-dense}
For \(F\ge 3\), the set \(U_{\mathrm{label}}\) contains a Zariski-open dense
subset of \((V^*)^F\). In particular, \(U_{\mathrm{label}}\) is dense in the
Euclidean topology. It is also Euclidean open.
\end{lemma}

\begin{proof}
For each pair \(r\ne s\), consider the condition
\[
  p_r(A)\ne p_s(A).
\]
It is the union, over \(1\le k\le F\), of the polynomial nonvanishing
conditions \(p_{r,k}(A)-p_{s,k}(A)\ne 0\). We first check that this union is not
empty. Choose \(a_r=e_1\), \(a_s=e_2\), and one further row \(a_t=e_3\). Put all
other rows equal to \(0\). Then \(\omega(e_1,e_3)=1\) and
\(\omega(e_2,e_3)=0\), so \(p_{r,2}(A)\ne p_{s,2}(A)\). Therefore, for every
pair \(r\ne s\), the polynomial differences \(p_{r,k}-p_{s,k}\) are not all
identically zero.

The intersection over all pairs \(r\ne s\) of these nonvanishing unions is
therefore Zariski open dense. On this intersection, the power sums of the row
multisets differ for every pair of rows, so the sorted row signatures
\(\sigma_i(A)\) are pairwise distinct. Hence this Zariski-open dense set is
contained in \(U_{\mathrm{label}}\).

Finally, \(U_{\mathrm{label}}\) is Euclidean open because the finitely many
sorted tuples vary continuously as unordered multisets, and pairwise distinct
compact tuples have a positive separation.
\end{proof}

\begin{definition}[Symplectic Gram--Schmidt for an ordered quadruple]
\label{def:symplectic-gram-schmidt-quadruple}
For an ordered quadruple \((u_1,u_2,u_3,u_4)\), set
\[
  q_1=u_1.
\]
If \(\omega(q_1,u_2)\ne 0\), set
\[
  p_1=\frac{u_2}{\omega(q_1,u_2)}.
\]
Define
\[
  \Pi(v)=v-\omega(v,p_1)q_1+\omega(v,q_1)p_1.
\]
Set
\[
  q_2=\Pi(u_3).
\]
If \(\omega(q_2,\Pi(u_4))\ne 0\), set
\[
  p_2=\frac{\Pi(u_4)}{\omega(q_2,\Pi(u_4))}.
\]
When both denominators are nonzero and \(q_2\ne 0\), the quadruple succeeds and
produces
\[
  F=[q_1\ q_2\ p_1\ p_2].
\]
\end{definition}

\begin{definition}[Generic coordinate-section domain and representative]
\label{def:generic-coordinate-section}
Let \(U\subset (V^*)^F\) be the set of row lists satisfying:
\begin{enumerate}
  \item \(A\in U_{\mathrm{label}}\);
  \item \(a_1,\ldots,a_F\) span \(V^*\);
  \item the lexicographic scan of ordered quadruples from \(B_A\) has a first
    successful quadruple under
    Definition~\ref{def:symplectic-gram-schmidt-quadruple}.
\end{enumerate}
For \(A\in U\), let \(F_A\) be the frame from that first successful quadruple.
Define
\[
  C(A)=\bigl(F_A^{-1}b_1,\ldots,F_A^{-1}b_F\bigr).
\]
This is a coordinate row list, not an invariant matrix.
\end{definition}

\begin{lemma}[The coordinate-section domain is open dense]
\label{lem:generic-coordinate-section-open-dense}
For \(F\ge 4\), the set \(U\) from
Definition~\ref{def:generic-coordinate-section} is Euclidean open dense in
\((V^*)^F\).
\end{lemma}

\begin{proof}
The labeling condition is open dense by
Lemma~\ref{lem:generic-coordinate-labeling-open-dense}. The spanning condition
is also open dense for \(F\ge 4\): its complement is defined by the vanishing of
all \(4\times4\) minors of the \(F\times4\) row matrix, and at least one such
minor is not identically zero.

It remains to check that a spanning row list has at least one successful
ordered quadruple. Choose four rows that form a basis of \(V^*\). The
restriction of \(\omega\) to their span is nondegenerate because their span is
all of \(V^*\). Hence among the four rows there is an ordered pair with
nonzero \(\omega\)-pairing. After the first hyperbolic pair is chosen, the
symplectic orthogonal quotient has dimension two and is nondegenerate. The
projections of the remaining basis rows span that quotient, so some ordering of
them has nonzero \(\omega\)-pairing after projection. Therefore the exhaustive
ordered-quadruple scan succeeds.

Thus, on the spanning locus, the third item in
Definition~\ref{def:generic-coordinate-section} is automatic. The set \(U\) is
therefore the intersection of the open dense labeling and spanning conditions.
\end{proof}

\begin{lemma}[The frame is symplectic]
\label{lem:generic-coordinate-frame-symplectic}
For \(A\in U\), the matrix \(F_A\) satisfies
\[
  F_A^TJF_A=J.
\]
\end{lemma}

\begin{proof}
By construction, \(\omega(q_1,p_1)=1\). The map \(\Pi\) subtracts the
\(\langle q_1,p_1\rangle\) components, so
\[
  \omega(\Pi(v),q_1)=\omega(\Pi(v),p_1)=0
\]
for every \(v\). Hence \(q_2=\Pi(u_3)\) and \(\Pi(u_4)\) are symplectically
orthogonal to \(q_1\) and \(p_1\). The normalization defining \(p_2\) gives
\(\omega(q_2,p_2)=1\). Skew-symmetry gives the remaining diagonal zero entries,
and the orthogonality just proved gives the off-block zero entries. In the
ordered basis \((q_1,q_2,p_1,p_2)\), the matrix of \(\omega\) is \(J\).
\end{proof}

\begin{lemma}[Permutation equivariance of the row labeling]
\label{lem:generic-coordinate-labeling-permutation-equivariant}
For \(A\in U_{\mathrm{label}}\) and any permutation \(\tau\in S_F\), the ordered
list \(B_{\tau A}\) equals \(B_A\).
\end{lemma}

\begin{proof}
The matrix \(\Omega(\tau A)\) is obtained from \(\Omega(A)\) by simultaneous
row-column permutation. Therefore the row signature attached to a facet moves
with that facet. Since the signatures are pairwise distinct, sorting them
recovers the same ordered list of covectors.
\end{proof}

\begin{lemma}[Symplectic equivariance of the selected frame]
\label{lem:generic-coordinate-frame-symplectic-equivariant}
Let \(R\in Sp(4)\), and let
\[
  R^{-T}A=(R^{-T}a_1,\ldots,R^{-T}a_F).
\]
If \(A\in U\), then \(R^{-T}A\in U\), \(B_{R^{-T}A}=R^{-T}B_A\), and
\[
  F_{R^{-T}A}=R^{-T}F_A.
\]
\end{lemma}

\begin{proof}
For all \(i,j\),
\[
  \omega(R^{-T}a_i,R^{-T}a_j)
  =a_i^T R^{-1}JR^{-T}a_j
  =a_i^TJa_j
  =\omega(a_i,a_j).
\]
Thus the row signatures are unchanged and the same ordered labels are selected.
The lexicographic quadruple scan is therefore the same scan applied to the
transformed ordered list. Symplectic Gram--Schmidt commutes with \(R^{-T}\)
because its formulas use vector addition, scalar multiplication, and
\(\omega\)-pairings, all of which are preserved by \(R^{-T}\). Hence the first
successful quadruple is the transformed first successful quadruple and
\(F_{R^{-T}A}=R^{-T}F_A\).
\end{proof}

\begin{proposition}[Generic coordinate canonization]
\label{prop:generic-coordinate-canonization}
For \(A\in U\), the representative \(C(A)\) is invariant under
\[
  Sp(4)\times S_F.
\]
Equivalently, the construction gives a generic coordinate representative for
the symplectic and facet-permutation action.
\end{proposition}

\begin{proof}
Permutation invariance follows from
Lemma~\ref{lem:generic-coordinate-labeling-permutation-equivariant}: the ordered
list \(B_A\) is unchanged by input permutation, so the selected frame and
coordinate rows are unchanged.

For \(R\in Sp(4)\), Lemma~\ref{lem:generic-coordinate-frame-symplectic-equivariant}
gives \(B_{R^{-T}A}=R^{-T}B_A\) and \(F_{R^{-T}A}=R^{-T}F_A\). For each ordered
row \(b_i\),
\[
  F_{R^{-T}A}^{-1}R^{-T}b_i
  =(R^{-T}F_A)^{-1}R^{-T}b_i
  =F_A^{-1}b_i.
\]
Thus \(C(R^{-T}A)=C(A)\).
\end{proof}

\begin{remark}[Generic, not universal]
\label{rem:generic-coordinate-canonization-not-universal}
The current construction is not a universal section. It depends on the generic
condition that omega row signatures are distinct. Symmetric or near-symmetric
inputs can have signature ties, and then signature sorting is not a canonical
labeling algorithm. A universal version must replace
Definition~\ref{def:generic-row-labeling-domain} by a full canonical-labeling
or tie-backtracking rule.
\end{remark}
